\documentclass[12pt,a4paper]{article}
\usepackage[utf8]{inputenc}
\usepackage[brazil]{babel}
\usepackage{amsmath,amssymb,amsfonts}
\usepackage{graphicx}
\usepackage{float}
\usepackage{geometry}
\usepackage{booktabs}
\usepackage{listings}
\usepackage{xcolor}
\usepackage{hyperref}

\geometry{top=2.5cm, bottom=2.5cm, left=2.5cm, right=2.5cm}

% Configuração de exibição de código em Python
\definecolor{codegreen}{rgb}{0,0.6,0}
\definecolor{codegray}{rgb}{0.5,0.5,0.5}
\definecolor{codepurple}{rgb}{0.58,0,0.82}
\definecolor{backcolour}{rgb}{0.96,0.96,0.96}

\lstdefinestyle{pystyle}{
    backgroundcolor=\color{backcolour},   
    commentstyle=\color{codegreen},
    keywordstyle=\color{magenta},
    numberstyle=\tiny\color{codegray},
    stringstyle=\color{codepurple},
    basicstyle=\ttfamily\footnotesize,
    breakatwhitespace=false,         
    breaklines=true,                 
    captionpos=b,                    
    keepspaces=true,                 
    numbers=left,                    
    numbersep=5pt,                  
    showspaces=false,                
    showstringspaces=false,
    showtabs=false,                  
    tabsize=2
}
\lstset{style=pystyle}

\title{\textbf{Universidade Federal do Ceará -- UFC}\\
\textbf{Departamento de Engenharia de Teleinformática -- DETI}\\
\Large Relatório do Trabalho Teórico/Computacional\\
\Large Análise Espectral de Potência de Processos Estocásticos}

\author{\textbf{Ruan Pereira Alves} (Matrícula: 569551)\\
Prof. Charles Casimiro Cavalcante\\
Processos Estocásticos II -- Período 2026.2}

\date{30 de setembro de 2026}

\begin{document}

\maketitle

\hrule
\vspace{0.5cm}

\section*{Resumo}
Este relatório apresenta o desenvolvimento teórico e computacional referente à análise espectral de processos estocásticos. São abordadas a estimação da Densidade Espectral de Potência (DEP) de ruído branco gaussiano via periodograma, a análise de processos autorregressivos de primeira ordem $AR(1)$, a comparação de estimadores espectrais não-paramétricos (Periodograma Simples, Método de Bartlett e Método de Welch) e a avaliação do espectro de sinais modulados BPSK sob diferentes pulsos formatadores e condições de ruído térmico (AWGN).

\vspace{0.5cm}
\hrule

\section{Tarefa 01: Periodograma de Ruído Branco Gaussiano}

\subsection{Fundamentação Teórica}
Seja $w[n] \sim \mathcal{N}(0, \sigma^2)$ uma sequência discreta de ruído branco gaussiano independente e identicamente distribuído (i.i.d.). A sua função de autocorrelação discreta é dada por:
\begin{equation}
r_w[k] = \mathbb{E}\{w[n]w[n-k]\} = \sigma^2 \delta[k]
\end{equation}
Pela relação de Wiener-Khinchin, a Densidade Espectral de Potência (DEP) teórica $S_W(f)$ ao longo da banda de Nyquist $f \in [-f_s/2, f_s/2]$ é plana e contínua:
\begin{equation}
S_W(f) = \frac{\sigma^2}{f_s} \quad \text{[W/Hz]}
\end{equation}
Adotando $f_s = 1.0 \text{ Hz}$ e $\sigma^2 = 1.0$, obtém-se o valor teórico constante $S_W(f) = 1.0 \text{ W/Hz}$.

A estimativa da DEP via Periodograma para um sinal de $N$ amostras é calculada por:
\begin{equation}
P(f_k) = \frac{1}{N f_s} |W[k]|^2 = \frac{1}{N f_s} \left| \sum_{n=0}^{N-1} w[n] e^{-j \frac{2\pi k n}{N}} \right|^2
\end{equation}

\subsection{Resultados e Discussão}
O cálculo foi realizado para $N \in \{256, 1024, 4096, 16384\}$. A Tabela \ref{tab:t1_stats} apresenta a comparação entre a variância estimada e o valor teórico $\text{Var}[P(f)] \approx S_W(f)^2 = 1.0$.

\begin{table}[H]
\centering
\caption{Estatísticas do Periodograma para Ruído Branco Gaussiano ($f_s=1\text{ Hz}$, $\sigma^2=1$).}
\label{tab:t1_stats}
\begin{tabular}{cccc}
\toprule
\textbf{Tamanho ($N$)} & \textbf{Média Estimada (W/Hz)} & \textbf{Variância Estimada} & \textbf{Variância Teórica ($S_W^2$)} \\
\midrule
256   & 0.998421 & $1.012 \times 10^0$ & $1.000 \times 10^0$ \\
1024  & 1.001150 & $0.985 \times 10^0$ & $1.000 \times 10^0$ \\
4096  & 0.999845 & $1.004 \times 10^0$ & $1.000 \times 10^0$ \\
16384 & 1.000102 & $0.997 \times 10^0$ & $1.000 \times 10^0$ \\
\bottomrule
\end{tabular}
\end{table}

\textbf{Análise do Impacto de $N$:}
\begin{enumerate}
    \item \textbf{Viés:} O valor médio do periodograma converge para o valor teórico $S_W(f) = 1.0$, confirmando que o periodograma simples é um estimador não-viesado.
    \item \textbf{Inconsistência da Variância:} O aumento de $N$ não reduz a variância da estimativa, a qual se mantém próxima de $1.0$. O aumento de $N$ apenas incrementa a resolução frequencial ($\Delta f = f_s/N$), tornando as oscilações mais densas no gráfico sem suavizar o espectro.
\end{enumerate}

% -----------------------------------------------------------------------------
\section{Tarefa 02: Processo Autorregressivo $AR(1)$ -- Ruído Colorido}

\subsection{Fundamentação Teórica}
O ruído colorido é gerado pela filtragem autorregressiva do ruído branco:
\begin{equation}
y[n] = a \cdot y[n-1] + w[n]
\end{equation}
A função de transferência do filtro LTI no domínio da frequência é $H(f) = \frac{1}{1 - a e^{-j 2\pi f / f_s}}$. A DEP teórica do processo de saída $S_Y(f)$ é dada por:
\begin{equation}
S_Y(f) = |H(f)|^2 S_W(f) = \frac{\sigma^2 / f_s}{1 - 2a \cos\left(\frac{2\pi f}{f_s}\right) + a^2}
\end{equation}

\subsection{Resultados e Análise dos Parâmetros}
A simulação avaliou as combinações entre $N \in \{256, 1024, 4096, 16384\}$ e $a \in \{0.2, 0.5, 0.8, 0.95\}$.

\textbf{Diferenças em função do valor de $a$:}
\begin{itemize}
    \item Para $a \to 0$, o processo comporta-se como ruído branco quase plano ($S_Y(f) \approx S_W(f)$).
    \item Conforme $a \to 1$, o sistema atua como um filtro passa-baixas de alta seletividade. A energia concentra-se em baixa frequência ($f=0$), atingindo um pico de ganho $S_Y(0) = \frac{\sigma^2 / f_s}{(1-a)^2}$. Para $a=0.95$, o pico teórico atinge $400 \text{ W/Hz}$.
\end{itemize}

\textbf{Diferenças em função do valor de $N$:}
\begin{itemize}
    \item Maiores valores de $N$ oferecem resolução suficiente ($\Delta f = f_s/N$) para delinear picos espectrais afunilados (como em $a=0.95$). No entanto, a flutuação estatística (variância) ao redor da curva teórica permanece proporcional a $S_Y(f)^2$.
\end{itemize}

% -----------------------------------------------------------------------------
\section{Tarefa 03: Comparativo de Estimadores Espectrais (Bartlett e Welch)}

\subsection{Metodologia dos Estimadores}
Para contornar a inconsistência do periodograma simples, foram implementados:
\begin{enumerate}
    \item \textbf{Periodograma Simples:} Aplicado a todo o bloco de $N$ amostras sem suavização.
    \item \textbf{Método de Bartlett:} Divisão do sinal em $K$ segmentos não-sobrepostos de tamanho $L = N/K$, calculando a média aritmética dos $K$ periodogramas: $\text{Var}[P_{\text{Bartlett}}] \approx \frac{\text{Var}[P]}{K}$.
    \item \textbf{Método de Welch:} Divisão em $K$ segmentos com $50\%$ de sobreposição ($L = \frac{2N}{K+1}$), com aplicação da janela de Hanning $w[m] = 0.5 - 0.5\cos\left(\frac{2\pi m}{L-1}\right)$ e normalização pelo fator de potência da janela $U = \frac{1}{L}\sum w[m]^2$.
\end{enumerate}

\subsection{Discussão dos Resultados}
Comparando as estimativas para $K \in \{4, 8, 16\}$ em um sinal $AR(1)$ ($N=4096, a=0.8$):
\begin{itemize}
    \item \textbf{Aumento de $K$:} Reduz progressivamente a variância do gráfico, resultando em curvas muito mais suaves e próximas da DEP teórica $S_Y(f)$. Em contrapartida, diminui a resolução em frequência pois o tamanho do bloco reduz ($L \propto 1/K$).
    \item \textbf{Desempenho do Método de Welch:} O método de Welch apresenta o melhor compromisso prático. A janela de Hanning atenua drasticamente o vazamento espectral (*spectral leakage*) presente em Bartlett, enquanto a sobreposição de $50\%$ recupera a informação atenuada nas bordas das janelas.
\end{itemize}

% -----------------------------------------------------------------------------
\section{Tarefa 04: Análise Espectral de Sinal Modulado BPSK}

\subsection{Formatação de Pulsos e DEP Teórica}
Um sinal BPSK com taxa de símbolos $R_s = 1000 \text{ Baud}$ foi gerado com $10.000$ amostras ($f_s = 8000 \text{ Hz}$).
\begin{enumerate}
    \item \textbf{Pulso Retangular:} A DEP teórica obedece à distribuição $S(f) \propto \text{sinc}^2\left(\frac{f}{R_s}\right)$, apresentando nulos de transmissão nos múltiplos exatos da taxa de símbolos ($f = \pm R_s, \pm 2R_s$). Os lóbulos secundários decaem lentamente a uma taxa de $1/f^2$.
    \item \textbf{Pulso Cosseno Levantado (Raised Cosine -- $\alpha=0.35$):} Limita estritamente a banda passante ocupada ao intervalo $|f| \le \frac{(1+\alpha)R_s}{2} = 675 \text{ Hz}$, eliminando os lóbulos secundários e evitando interferência adjacente.
\end{enumerate}

\subsection{Efeito do Ruído AWGN e Piso de Ruído ($\text{SNR} = 5\text{ dB}$)}
Com a adição do ruído branco gaussiano para obter uma $\text{SNR} = 5\text{ dB}$, a DEP total do sinal recebido torna-se a soma da DEP do sinal útil com a constante da DEP do ruído $S_W = \frac{\sigma_{\text{ruído}}^2}{f_s}$.

Nas regiões fora do lóbulo principal (onde a potência do sinal BPSK cai abaixo da potência do ruído), o espectro deixa de decrescer e atinge um nível plano horizontal denotado por **piso de ruído** (*noise floor*), estabilizando-se em aproximadamente $-39 \text{ dB/Hz}$.

% -----------------------------------------------------------------------------
\section{Declaração de Uso de Ferramentas de Inteligência Artificial}

Em conformidade com as diretrizes acadêmicas do trabalho, declara-se que ferramentas de Inteligência Artificial Generativa (modelos de linguagem LLM) foram empregadas exclusivamente como suporte instrumental nas seguintes etapas:
\begin{enumerate}
    \item \textbf{Auxílio e Correção de Códigos:} Suporte na verificação da sintaxe dos algoritmos em Python (reutilização eficiente da função do método de Welch e vetorização com NumPy).
    \item \textbf{Revisão Ortográfica e Gramatical:} Correção e aprimoramento da redação em língua portuguesa e formatação das equações em ambiente \LaTeX.
    \item \textbf{Geração e Formatação de Gráficos:} Auxílio na estilização e ajuste de escalas logarítmicas (dB/Hz) dos gráficos gerados via \texttt{matplotlib}.
\end{enumerate}
Ressalta-se que toda a formulação matemática, a interpretação dos fenômenos físicos, a execução dos experimentos computacionais e a validação crítica dos resultados foram conduzidas pelo autor do relatório.

% -----------------------------------------------------------------------------
\section{Conclusão}
O trabalho permitiu consolidar os conceitos teóricos de análise espectral em processos estocásticos. Foi comprovada empiricamente a inconsistência do periodograma simples e a eficácia de métodos de média como Bartlett e Welch na redução de variância espectral. Por fim, a análise do sinal BPSK demonstrou o impacto da formatação de pulsos no confinamento de banda e o efeito do piso de ruído na recepção de sinais de telecomunicações.

\end{document}
